\documentclass[12pt]{article}
\usepackage[utf8]{inputenc}
\usepackage[spanish]{babel}
\usepackage{geometry}
\usepackage{xcolor}
\usepackage{listings}
\usepackage{hyperref}
\usepackage{graphicx}

% Configuración de la página
\geometry{a4paper, left=2.5cm, right=2.5cm, top=2.5cm, bottom=2.5cm}

% Colores para el código
\definecolor{codegreen}{rgb}{0,0.6,0}
\definecolor{codegray}{rgb}{0.5,0.5,0.5}
\definecolor{codepurple}{rgb}{0.58,0,0.82}
\definecolor{backcolour}{rgb}{0.95,0.95,0.92}

% Estilo para los listados de código de Python
\lstdefinestyle{mystyle}{
    backgroundcolor=\color{backcolour},   
    commentstyle=\color{codegreen},
    keywordstyle=\color{magenta},
    numberstyle=\tiny\color{codegray},
    stringstyle=\color{codepurple},
    basicstyle=\ttfamily\footnotesize,
    breakatwhitespace=false,         
    breaklines=true,                 
    captionpos=b,                    
    keepspaces=true,                 
    numbers=left,                    
    numbersep=5pt,                  
    showspaces=false,                
    showstringspaces=false,
    showtabs=false,                  
    tabsize=2
}
\lstset{style=mystyle, language=Python}

% Título del documento
\title{Informe de Desarrollo con Asistencia de IA: \\ Análisis Bibliométrico con Grafos}
\author{Tu Nombre}
\date{\today}

\begin{document}

\maketitle

\begin{abstract}
Este documento detalla el proceso de desarrollo de dos scripts en Python, \texttt{requerimiento\_1\_grafos.py} y \texttt{requerimiento\_2\_grafos.py}, para un proyecto de análisis bibliométrico. Se reconstruye la secuencia de interacciones con un asistente de inteligencia artificial (IA) para ilustrar cómo fue utilizado como una herramienta de apoyo conceptual y técnico. El objetivo es demostrar un flujo de trabajo colaborativo donde la IA proporciona guía y código base, mientras que el estudiante es responsable de la integración, adaptación y desarrollo final.
\end{abstract}

\tableofcontents
\newpage

\section{Requerimiento 1: Grafo de Citaciones}
\label{sec:req1}

El objetivo de este requerimiento era construir y analizar una red de citaciones a partir de una colección de artículos en formato \texttt{.bib}. El grafo debía modelar tanto las citas explícitas como las inferidas por similitud de texto.

\subsection{Interacción 1: Creación del Grafo Básico}

\subsubsection*{Prompt (Usuario)}
\begin{quote}
Hola, necesito crear un grafo de citaciones a partir de un archivo \texttt{.bib}. Cada artículo es un nodo. Debería poder leer el archivo y, si un artículo A cita a un artículo B, se debe crear una arista dirigida de A a B. Por ahora, asumamos que la información de la cita está en un campo llamado \texttt{cites} dentro de cada entrada del \texttt{.bib}.
\end{quote}

\subsubsection*{Respuesta (IA)}
\begin{quote}
¡Claro! Para esa tarea, te recomiendo usar \texttt{bibtexparser} para leer los datos y \texttt{networkx} para el grafo. El plan sería:
\begin{enumerate}
    \item Define una función que parsee el archivo \texttt{.bib}.
    \item Crea un grafo dirigido (\texttt{nx.DiGraph}).
    \item Itera sobre cada artículo, añadiéndolo como un nodo.
    \item Si el artículo tiene el campo \texttt{cites}, procesa esa cadena y añade las aristas correspondientes.
\end{enumerate}
Aquí tienes un código base para empezar:
\end{quote}
\begin{lstlisting}
import bibtexparser
import networkx as nx

def construir_grafo_base(bib_file_path):
    with open(bib_file_path, encoding='utf-8') as bibtex_file:
        bib_database = bibtexparser.load(bibtex_file)
    
    G = nx.DiGraph()
    articulos = bib_database.entries
    
    # Primero, anade todos los articulos como nodos
    for articulo in articulos:
        node_id = articulo.get('ID')
        if node_id:
            G.add_node(node_id, title=articulo.get('title', 'N/A'))
    
    # Ahora, itera de nuevo para crear las aristas
    for articulo in articulos:
        node_id = articulo.get('ID')
        if 'cites' in articulo:
            citados_ids = [c.strip() for c in articulo['cites'].split(',')]
            for citado_id in citados_ids:
                # Asegurate de que el nodo citado exista
                if G.has_node(citado_id):
                    G.add_edge(node_id, citado_id)
    return G
\end{lstlisting}

\subsection{Interacción 2: Inferencia de Citaciones por Similitud}

\subsubsection*{Prompt (Usuario)}
\begin{quote}
Gracias, funciona bien. Sin embargo, muchos de mis artículos no tienen el campo \texttt{cites}. ¿Podemos mejorar el script para que infiera una posible cita basándose en la similitud del texto entre el título y el abstract? El artículo más reciente debería citar al más antiguo.
\end{quote}

\subsubsection*{Respuesta (IA)}
\begin{quote}
Excelente idea. Para inferir relaciones, te sugiero usar \texttt{scikit-learn}. El proceso es: vectorizar los textos con \texttt{TfidfVectorizer}, calcular la \texttt{cosine\_similarity} y luego añadir aristas basadas en un umbral de similitud y la fecha de publicación.

Puedes empezar con una función como esta para obtener la matriz de similitud. Luego, puedes integrar la lógica de añadir las aristas en tu función principal.
\end{quote}
\begin{lstlisting}
from sklearn.feature_extraction.text import TfidfVectorizer
from sklearn.metrics.pairwise import cosine_similarity
import numpy as np

def calcular_matriz_similitud(articulos):
    # Concatena titulo y abstract
    textos = [art.get('title', '') + ' ' + art.get('abstract', '') for art in articulos]
    
    vectorizer = TfidfVectorizer(stop_words='english', min_df=2)
    tfidf_matrix = vectorizer.fit_transform(textos)
    
    # Calcula la similitud de todos contra todos
    matriz_similitud = cosine_similarity(tfidf_matrix)
    np.fill_diagonal(matriz_similitud, 0)
    
    # Ahora te queda iterar sobre esta matriz. Para cada par (i, j) 
    # con alta similitud, compara los anos de los articulos i y j 
    # para decidir la direccion de la arista.
    return matriz_similitud
\end{lstlisting}

\subsection{Interacción 3: Funciones de Análisis de Grafo}

\subsubsection*{Prompt (Usuario)}
\begin{quote}
Perfecto. Ahora que el grafo es más completo, quiero analizarlo. Necesito una función para encontrar el 'camino de influencia' (mayor similitud) entre dos artículos y otra para detectar 'comunidades' (componentes fuertemente conexas). ¿Qué funciones de \texttt{networkx} me sirven?
\end{quote}

\subsubsection*{Respuesta (IA)}
\begin{quote}
¡Buenas preguntas! \texttt{networkx} tiene herramientas perfectas para eso.
\begin{enumerate}
    \item Para el 'camino de influencia', usa \texttt{nx.dijkstra\_path()}. Como busca el menor peso, el truco es definir el peso como \texttt{1 - similitud}.
    \item Para las 'comunidades', la función es \texttt{nx.strongly\_connected\_components(G)}.
\end{enumerate}
Aquí tienes unas funciones de ejemplo que puedes adaptar:
\end{quote}
\begin{lstlisting}
import networkx as nx

def encontrar_camino_minimo(G, origen, destino):
    """Encuentra el camino de mayor similitud usando Dijkstra."""
    try:
        # Asume que tus aristas inferidas tienen un atributo 'weight'.
        path = nx.dijkstra_path(G, source=origen, target=destino, 
                               weight=lambda u, v, d: 1 - d.get('weight', 0))
        return path
    except nx.NetworkXNoPath:
        return f"No hay un camino entre {origen} y {destino}."

def encontrar_componentes_fuertemente_conexas(G):
    """Encuentra y devuelve las componentes fuertemente conexas."""
    return list(nx.strongly_connected_components(G))
\end{lstlisting}

\newpage

\section{Requerimiento 2: Grafo de Co-ocurrencia de Términos}
\label{sec:req2}

El segundo requerimiento consistía en identificar los temas principales de la colección de artículos. Para ello, se decidió construir un grafo de co-ocurrencia de términos clave extraídos de los abstracts.

\subsection{Interacción 1: Creación del Grafo de Co-ocurrencia}

\subsubsection*{Prompt (Usuario)}
\begin{quote}
Ahora quiero hacer un análisis de temas. ¿Puedes explicarme cómo crear un script que construya un grafo de co-ocurrencia de términos? Los nodos serían los términos y una arista entre ellos significaría que aparecen en el mismo abstract, con un peso que cuente las co-ocurrencias.
\end{quote}

\subsubsection*{Respuesta (IA)}
\begin{quote}
Claro, esa es una técnica estándar llamada análisis de co-ocurrencia. La lógica es iterar sobre cada abstract, extraer los términos y, para cada par de términos, añadir o reforzar una arista en un grafo no dirigido.

Aquí tienes un código base para esa lógica. Necesitarás una lista predefinida de \texttt{terminos\_clave} para que funcione.
\end{quote}
\begin{lstlisting}
import networkx as nx
import itertools
import re

def construir_grafo_coocurrencia(abstracts, terminos_clave):
    G = nx.Graph()
    G.add_nodes_from(terminos_clave)
    
    for abstract in abstracts:
        # Encuentra que terminos clave estan en este abstract
        terminos_en_abstract = [t for t in terminos_clave if re.search(r'\b' 
                               + re.escape(t) + r'\b', abstract, re.IGNORECASE)]
        
        # Crea aristas para cada par que co-ocurre
        for term1, term2 in itertools.combinations(terminos_en_abstract, 2):
            if G.has_edge(term1, term2):
                G[term1][term2]['weight'] += 1
            else:
                G.add_edge(term1, term2, weight=1)
    return G
\end{lstlisting}

\subsection{Interacción 2: Selección de Términos con TF-IDF}

\subsubsection*{Prompt (Usuario)}
\begin{quote}
El resultado con todas las palabras es muy ruidoso. ¿Cómo puedo usar TF-IDF para extraer solo los 50 términos más importantes (palabras o bigramas) y usar solo esos para el grafo?
\end{quote}

\subsubsection*{Respuesta (IA)}
\begin{quote}
Totalmente. Usar TF-IDF para pre-seleccionar los términos es el enfoque correcto. Debes vectorizar todo el corpus, sumar los scores de cada término en todos los documentos y quedarte con los mejores.

Esta función te ayudará a extraer esos términos clave. Luego puedes pasar la lista resultante a la función del grafo que hicimos antes.
\end{quote}
\begin{lstlisting}
from sklearn.feature_extraction.text import TfidfVectorizer
import numpy as np

def extraer_terminos_relevantes(abstracts, num_terminos=50):
    """Extrae terminos y bigramas relevantes usando TF-IDF."""
    vectorizer = TfidfVectorizer(stop_words='english', ngram_range=(1, 2), 
                                 max_df=0.85, min_df=3)
    tfidf_matrix = vectorizer.fit_transform(abstracts)
    
    sum_tfidf = tfidf_matrix.sum(axis=0)
    scores = np.array(sum_tfidf).flatten()
    feature_names = np.array(vectorizer.get_feature_names_out())
    
    sorted_indices = np.argsort(scores)[::-1]
    mejores_terminos = feature_names[sorted_indices][:num_terminos]
    
    return list(mejores_terminos)
\end{lstlisting}

\subsection{Interacción 3: Filtrado y Análisis del Grafo Final}

\subsubsection*{Prompt (Usuario)}
\begin{quote}
Mucho mejor. ¿Y cómo puedo ahora filtrar las conexiones débiles del grafo y luego analizarlo para encontrar los términos más centrales y los clústeres temáticos?
\end{quote}

\subsubsection*{Respuesta (IA)}
\begin{quote}
Estás en el último paso. Primero, elimina las aristas con peso bajo y los nodos que queden aislados. Después, \texttt{networkx} te facilita el análisis:
\begin{itemize}
    \item \textbf{Centralidad:} La mides con el grado de los nodos (\texttt{G.degree()}).
    \item \textbf{Clústeres:} Los obtienes con las componentes conexas (\texttt{nx.connected\_components(G)}).
\end{itemize}
Aquí tienes un ejemplo de cómo podrías estructurar ese análisis final:
\end{quote}
\begin{lstlisting}
import networkx as nx

def filtrar_y_analizar_grafo(G, umbral_peso=5):
    # 1. Filtrar aristas debiles
    aristas_a_eliminar = [(u, v) for u, v, data in G.edges(data=True) 
                          if data['weight'] < umbral_peso]
    G.remove_edges_from(aristas_a_eliminar)
    
    # 2. Eliminar nodos aislados
    G.remove_nodes_from(list(nx.isolates(G)))
    print(f"Grafo filtrado con {G.number_of_nodes()} nodos.")

    # 3. Analizar centralidad
    grados = dict(G.degree())
    terminos_centrales = sorted(grados.items(), key=lambda item: item[1], reverse=True)
    print("\n--- Terminos mas conectados ---")
    print(terminos_centrales[:10])
    
    # 4. Analizar clusteres
    componentes = list(nx.connected_components(G))
    print("\n--- Grupos Tematicos ---")
    for i, comp in enumerate(componentes):
        if len(comp) > 2:
            print(f"Tema {i+1}: {list(comp)}")
\end{lstlisting}

\end{document}
